\documentclass[10pt,a4paper]{amsart}
\usepackage[margin=19mm]{geometry}
\usepackage{amsmath,amssymb,mathtools}
\usepackage[scr=rsfs,cal=euler]{mathalfa}
\usepackage{xfrac}
\usepackage[unicode,hidelinks,bookmarks=false]{hyperref}
\newcommand{\ie}{i.e.\ }
\newcommand{\cf}{cf.\ }
\newcommand{\resp}{respectively\ }
\newcommand{\Subsection}{\subsection}
\providecommand{\nobreakdash}{\nobreak\hbox{-}}
\setlength{\emergencystretch}{2em}
\multlinegap0pt
\usepackage{amssymb}
\usepackage{mathtools}
\newtheorem{lem}{Lemma}
\newtheorem{thm}{Theorem}
\newcommand{\C}{\mathbb C}
\title{Replication Descent and $J$-Finality of Replicable Functions}
\author{Eric Culf, Abdellah Sebbar}
\date{Source: arXiv:2609.08243v1 (CC BY 4.0)}
\begin{document}
\maketitle

\noindent\emph{Source and licence.} Replication Descent and $J$-Finality of Replicable Functions. Version arXiv:2609.08243v1 (CC BY 4.0), distributed by arXiv under the Creative Commons Attribution 4.0 International licence, http://creativecommons.org/licenses/by/4.0/, as recorded in the arXiv licence metadata for this exact version and shown on the abstract page https://arxiv.org/abs/2609.08243v1. Full licence notice: \texttt{licenses/arxiv-2609.08243-CC-BY-4.0.txt}.\par\smallskip

\noindent\textit{Editorial scope.} Counted unit: the article's thm thm:congruence-descent (source0.tex lines 335--341). The article's complete original proof of it is delivered with it (source0.tex lines 343--449, one \texttt{proof} environment of 107 lines). No mathematics is written, repaired or re-derived: the statement and the proof are the article's own lines, in the article's own notation, and no retained line is substituted or re-flowed.  Retained uncounted as the source-local context the proof needs: lines 268--273.  Nothing was omitted from any retained span, so the package carries no omission record.  No retained line contains a citation, so the wrapper carries no bibliography and no bibliography entry.  No retained line contains a figure environment, an image-inclusion command or drawing code, so no image is delivered and the package declares no asset dependency.  Editorial formatting only, in addition: an AMS \texttt{amsart} preamble that restates the article's own declarations and macros used by the retained lines, namely the environments \texttt{lem}, \texttt{thm} on separate counters, together with the standard packages those lines need.  The retained environments print in this document as Lemma 1, Theorem 1 in order of appearance, whereas the article prints the same items with the numbering of its own full document.  The complete native source of the article as distributed in the arXiv e-print for this exact version is bundled unchanged as \texttt{sources/209705/source0.tex}.
\begin{lem}\label{lem:generation}
For all positive integers $m,n$,
\begin{equation}\label{eq:generation}
\Gamma_0(m)=\left\langle\Gamma_0(m,n),T\right\rangle.
\end{equation}
\end{lem}

\begin{thm}\label{thm:congruence-descent}
Let $f$ be a normalized replicable function. Suppose that $f$ is invariant under $\Gamma_0(N)$. Then for every $n\geq1$,
\begin{equation}\label{eq:descent-level}
f^{(n)}\ \text{is invariant under}\ 
\Gamma_0\!\left(\frac{N}{(N,n)}\right).
\end{equation}
\end{thm}

\begin{proof}
The proof is by induction on $n$. For $n=1$ the assertion is the hypothesis. Assume it holds for every proper divisor of $n$ and set
\[
L=\operatorname{lcm}(N,n),
\qquad
M_a=\frac{N}{(N,a)}
\quad(a\mid n).
\]
One has the divisibility
\begin{equation}\label{eq:lcm-divisibility}
M_a=\frac{N}{(N,a)}\mid\frac{L}{a}.
\end{equation}
Indeed,
\[
aM_a=\operatorname{lcm}(a,N)
\]
divides $\operatorname{lcm}(n,N)=L$ because $a\mid n$.

Fix a factorization $ad=n$ with $a<n$ and define
\begin{equation}\label{eq:branch-sum}
S_{a,d}(\tau)=
\sum_{b\bmod d}
 f^{(a)}\!\left(\frac{a\tau+b}{d}\right).
\end{equation}
By the induction hypothesis, $f^{(a)}$ is invariant under $\Gamma_0(M_a)$. It remains to show that $S_{a,d}$ is invariant under $\Gamma_0(L)$.

Take
\[
\gamma=\begin{pmatrix}A&B\\C&D\end{pmatrix}\in\Gamma_0(L).
\]
Since $d\mid n\mid L$, reduction of $AD-BC=1$ modulo $d$ gives $AD\equiv1\pmod d$, so $A$ is invertible modulo $d$. For every $b\bmod d$ there is therefore a unique $b'\bmod d$ satisfying
\begin{equation}\label{eq:bprime}
Ab'\equiv aB+bD\pmod d.
\end{equation}
Consider
\[
\delta_{b}=\alpha_{a,b,d}\,\gamma\,\alpha_{a,b',d}^{-1}.
\]
A direct multiplication gives
\begin{equation}\label{eq:delta-matrix}
\delta_b=
\begin{pmatrix}
A+\dfrac{bC}{a}
&
\dfrac{aB+bD-Ab'}{d}-\dfrac{bb'C}{ad}
\\[3mm]
\dfrac{dC}{a}
&
D-\dfrac{b'C}{a}
\end{pmatrix}.
\end{equation}
All four entries are integers. For the diagonal entries this follows from $a\mid n\mid L\mid C$. The first term in the upper-right entry is integral by \eqref{eq:bprime}, and the second is integral because $ad=n\mid C$. The determinant is one because $\delta_b$ is a conjugated product of determinant-one.

It remains to check the lower-left congruence. By \eqref{eq:lcm-divisibility}, $M_a\mid L/a$, while $L\mid C$. Hence
\[
M_a\mid\frac{C}{a}\mid\frac{dC}{a}.
\]
Thus $\delta_b\in\Gamma_0(M_a)$. Consequently
\[
f^{(a)}\bigl(\alpha_{a,b,d}\gamma\tau\bigr)
=
f^{(a)}\bigl(\delta_b\alpha_{a,b',d}\tau\bigr)
=
f^{(a)}\bigl(\alpha_{a,b',d}\tau\bigr).
\]
The map $b\mapsto b'$ is a permutation modulo $d$: by \eqref{eq:bprime} it is an affine map with linear coefficient $A^{-1}D$, a unit modulo $d$. Thus summing over $b$ proves
\[
S_{a,d}(\gamma\tau)=S_{a,d}(\tau).
\]
Return to the replication identity and isolate the term $a=n,d=1$:
\begin{equation}\label{eq:isolate-n-descent}
f^{(n)}(n\tau)
=
\Phi_n(f(\tau))
-
\sum_{\substack{ad=n\\a<n}}S_{a,d}(\tau).
\end{equation}
Because $N\mid L$, the function $\Phi_n(f)$ is invariant under $\Gamma_0(L)$. Every branch sum on the right is also invariant under $\Gamma_0(L)$ by the preceding argument. Hence
\[
g_n(\tau):=f^{(n)}(n\tau)
\]
is invariant under $\Gamma_0(L)$.

Let
\[
\beta_n=\begin{pmatrix}n&0\\0&1\end{pmatrix}.
\]
Since $g_n=f^{(n)}\circ\beta_n$, conjugation gives invariance of $f^{(n)}$ under
\begin{equation}\label{eq:conjugate-subgroup}
\beta_n\Gamma_0(L)\beta_n^{-1}
=
\Gamma_0(L/n)\cap\Gamma^0(n).
\end{equation}
Here
\[
\frac{L}{n}=\frac{N}{(N,n)}.
\]
If $\gamma=\left(\begin{smallmatrix}A&B\\C&D\end{smallmatrix}\right)$ with $L\mid C$, then
\[
\beta_n\gamma\beta_n^{-1}
=
\begin{pmatrix}A&nB\\C/n&D\end{pmatrix},
\]
whose upper-right entry is divisible by $n$ and whose lower-left entry is divisible by $L/n$. Conversely every matrix in the intersection on the right of \eqref{eq:conjugate-subgroup} arises in this way.

The ordinary integral-power $q$-expansion of $f^{(n)}$ gives invariance under $T:\tau\mapsto\tau+1$. Lemma~\ref{lem:generation}, with $m=L/n=N/(N,n)$, then gives invariance under the full group $\Gamma_0(N/(N,n))$.
\end{proof}

\end{document}
